\documentclass[12pt,a4paper]{report}
\usepackage{graphicx}
\usepackage{hyperref}
\usepackage{listings}
\usepackage{xcolor}
\usepackage{geometry}
\usepackage{titlesec}
\usepackage{fancyhdr}
\usepackage{atbegshi}
\usepackage{background}
\usetikzlibrary{calc}
\usepackage{tikz}

% Define custom border style
\backgroundsetup{
  scale=1,
  angle=0,
  opacity=1,
  contents={
    \begin{tikzpicture}[remember picture, overlay]
      \draw [line width=1pt, black]
        ($(current page.north west) + (1cm,-1cm)$)
        rectangle
        ($(current page.south east) + (-1cm,1cm)$);
    \end{tikzpicture}
  }
}

% Add this line to apply the page border to every page
\geometry{a4paper, margin=1in}

% Title formatting
\titleformat{\chapter}[display]
{\normalfont\huge\bfseries}{\chaptertitlename\ \thechapter}{20pt}{\Huge}
\titlespacing*{\chapter}{0pt}{50pt}{40pt}

% Code listing style
\lstset{
    basicstyle=\ttfamily\small,
    breaklines=true,
    commentstyle=\color{green!50!black},
    keywordstyle=\color{blue},
    stringstyle=\color{red},
    numbers=left,
    numberstyle=\tiny\color{gray},
    numbersep=5pt,
    frame=single,
    framexleftmargin=15pt,
    framexrightmargin=0pt,
    framexbottommargin=5pt,
    framextopmargin=5pt,
    backgroundcolor=\color{gray!10},
    showstringspaces=false
}

\lstdefinelanguage{JavaScript}{
  keywords={typeof, new, true, false, catch, function, return, null, catch, switch, var, if, in, while, do, else, case, break},
  keywordstyle=\color{blue}\bfseries,
  ndkeywords={class, export, boolean, throw, implements, import, this},
  ndkeywordstyle=\color{darkgray}\bfseries,
  identifierstyle=\color{black},
  sensitive=false,
  comment=[l]{//},
  morecomment=[s]{/}{/},
  commentstyle=\color{purple}\ttfamily,
  stringstyle=\color{red}\ttfamily,
  morestring=[b]',
  morestring=[b]"
}

\renewcommand{\headrulewidth}{0.4pt}
\renewcommand{\footrulewidth}{0.4pt}

\pagestyle{fancy}
\fancyhf{}
\fancyhead[R]{E-Banking System}
\fancyfoot[C]{\thepage}

\begin{document}

\begin{titlepage}
    \centering
    \vspace*{1cm}
    {\huge\bfseries WT PROJECT REPORT\par}
    \vspace{1cm}
    {\Large\bfseries on\par}
    \vspace{1cm}
    {\huge\bfseries E-BANKING SYSTEM\par}
    \vspace{2cm}
    {\Large\bfseries Bachelor of Technology\par}
    \vspace{0.5cm}
    {\Large\bfseries In\par}
    \vspace{0.5cm}
    {\Large\bfseries Computer Science and Engineering\par}
    \vspace{2cm}
    {\Large\bfseries Under esteemed Guidance of\par}
    \vspace{0.5cm}
    {\Large\bfseries Mrs. Venkata Ratnam\par}
    \vspace{0.5cm}
    {\Large\bfseries Asst. Professor, Dept. of CSE\par}
    \vspace{2cm}
    {\Large\bfseries DEPARTMENT OF COMPUTER SCIENCE AND ENGINEERING\par}
    \vspace{0.5cm}
    {\Large\bfseries ANIL NEERUKONDA INSTITUTE OF TECHNOLOGY AND SCIENCES\par}
    \vspace{0.5cm}
    {\Large\bfseries SANGIVALASA, VISAKHAPATNAM\par}
\end{titlepage}

\thispagestyle{empty}
\begin{center}
    \LARGE\textbf{PROJECT TITLE: ONLINE BANKING SYSTEM}\\
    \vspace{2cm}

    \begin{tabular}{|l|l|}
        \hline
        \textbf{Name of the Student} & \textbf{Roll No} \\
        \hline
        D. Chaitanya & A22126510144 \\
        \hline
        P. Meghana & A22126510167 \\
        \hline
        P. Sai Deepak & A22126510168 \\
        \hline
        N. Lokesh & A22126510169 \\
        \hline
        M. Nikitha & A22126510161 \\
        \hline
    \end{tabular}
\end{center}
\newpage

\thispagestyle{empty}
\begin{center}
    \Large\textbf{Bonafide Certificate}
\end{center}
\vspace{1cm}

This is to certify that this project report "E-Banking System" is the Bonafede
work of D. Chaitanya(A22126510144), P. Meghana(A22126510167), P. Sai Deepak
(A22126510168), N. Lokesh(A22126510169), M. Nikitha(A22126510161). This project
is carried out and is submitted in the partial fulfillment of the requirements
for the award of BACHELOR OF TECHNOLOGY in Computer Science and Engineering,
under Anil Neerukonda Institute of Technology and Sciences during the academic
year 2024-2025.

\vspace{2cm}
\begin{tabular}{ll}
Head of the Department & Project Guide \\
Prof.M.Ramakrishna Murty & Mrs. K.Venkataratnam \\
(Professor \& HOD) & (Asst.Professor) \\
Department of CSE & Department of CSE \\
ANITS & ANITS
\end{tabular}
\newpage

\thispagestyle{empty}
\begin{center}
    \Large\textbf{Declaration}
\end{center}
\vspace{1cm}

This is to certify that the project work entitled "E-Banking System" is a
Bonafide work carried out by as a part of BTech 3\textsuperscript{rd} year
2\textsuperscript{nd} semester of Computer Science and Engineering of Anil
Neerukonda Institute of Technology and Sciences, Visakhapatnam during the
academic year 2024-2025.

We are D. Chaitanya, P. Meghana, P. Sai Deepak, N. Lokesh, M. Nikitha of 3rd
year B.Tech, Department of Computer Science and Engineering from ANITS,
Visakhapatnam, hereby declare that the project work entitled "E-Banking System"
is carried out by us and is submitted in the fulfillment of the requirement for
the award of Bachelor of Technology in Computer Science and Engineering, under
Anil Neerukonda Institute of Technology and Sciences during the academic year
2024-2025 has not been submitted to any other university for the award of any
kind of degree.

\vspace{2cm}
\begin{tabular}{ll}
\textbf{Name} & \textbf{Reg.No} \\
D. Chaitanya & A22126510144 \\
P. Meghana & A22126510167 \\
P. Sai Deepak & A22126510168 \\
N. Lokesh & A22126510169 \\
M. Nikitha & A22126510161 \\
\end{tabular}
\newpage

\thispagestyle{empty}
\begin{center}
    \Large\textbf{Acknowledgement}
\end{center}
\vspace{1cm}

The satisfaction that accompanies the successful completion of any task would
be incomplete without mentioning the people who made it possible, and whose
constant guidance and encouragement always upheld the morale. We take a great
pleasure in presenting a project, which is the result of a studied blend of
both research and knowledge.

We first take the privilege to thank our Head of the department Prof M.
RAMAKRISHNA MURTHY, for permitting us in laying the first stone of success. We
feel grateful to thank MRS. K. VENKATA RATNAM Madam, who is our project guide
and who shared her valuable knowledge with us and made us understand the real
essence of the topic and created interest in us to put our continuous efforts
in the project.

\vspace{2cm}
\begin{tabular}{ll}
\textbf{Name} & \textbf{Reg.No} \\
D. Chaitanya & A22126510144 \\
P. Meghana & A22126510167 \\
P. Sai Deepak & A22126510168 \\
Lokesh & A22126510169 \\
Nikitha & A22126510161 \\
\end{tabular}
\newpage

\tableofcontents
\newpage

\chapter{Project Abstract}

The proliferation of online banking has revolutionized the way financial transactions are
conducted, offering unparalleled convenience and accessibility. However, with this
convenience comes the challenge of ensuring robust security measures to protect sensitive
financial data and prevent unauthorized access.

This project aims to address these challenges by proposing innovative solutions to enhance
both security and user experience in online banking systems.

The project will begin with a comprehensive analysis of existing security protocols and user
interface designs in online banking systems. This analysis will identify potential
vulnerabilities and areas for improvement. Subsequently, the project will focus on the
development and implementation of advanced security measures such as multi-factor
authentication, encryption techniques, and real-time fraud detection algorithms.

Additionally, user-centric design principles will be employed to create intuitive and user-
friendly interfaces that enhance the overall banking experience.

The effectiveness of the proposed solutions will be evaluated through rigorous testing
procedures, including simulated cyber-attacks and user feedback surveys. Furthermore, the
project will explore the integration of emerging technologies such as biometrics and
blockchain to further enhance security and streamline transactions.

By the conclusion of this project, we anticipate significant advancements in the security and
usability of online banking systems. These advancements will not only bolster consumer
trust and confidence but also contribute to the continued growth and evolution of digital
banking services in an increasingly interconnected world.

\chapter{Objectives}

Here are a few objectives for your E-Banking System project:
\begin{enumerate}
    \item To develop a secure and user-friendly online banking platform that allows users to
    perform various banking transactions, such as account creation, fund transfers, and
    transaction history viewing.
    \item To implement robust security measures to protect user data and prevent unauthorized
    access to accounts.
    \item To provide an intuitive user interface that enhances the overall user experience and
    simplifies the process of managing finances online.
    \item To enable real-time monitoring of transactions and account activities for both users
    and administrators.
    \item To facilitate easy integration with existing banking systems and databases for
    seamless data management.
    \item To ensure compatibility with various devices and platforms, including web and
    mobile applications, to provide users with flexibility and convenience in accessing their
    accounts.
    \item To conduct thorough testing and validation of the system to ensure its reliability,
    performance, and security.
    \item To gather user feedback and make necessary improvements to enhance the system's
    functionality and usability.
    \item To explore the potential for future enhancements, such as the integration of
    emerging technologies like artificial intelligence and machine learning for fraud detection
    and personalized banking experiences.
\end{enumerate}


\chapter{Software \& Modules Used}

\begin{enumerate}
    \item \textbf{Deno:}
    Deno is a runtime for JavaScript and TypeScript, which is used to run the
    server-side code. It is a secure runtime that allows you to run JavaScript
    and TypeScript code outside of a web browser. It also serves a built-in
    package manager and supports ES modules.

    \item \textbf{Express.js:}
    Express.js is a web application framework for Node.js, designed for
    building web applications and APIs. It provides a robust set of features
    for web and mobile applications, including routing, middleware support, and
    template rendering.

    \item \textbf{MySQL:}
    MySQL is an open-source relational database management system. It is widely
    used for web applications and is known for its reliability and performance.
    MySQL allows you to store and retrieve data efficiently, and it supports
    SQL (Structured Query Language) for querying and managing databases. In
    this project, MySQL is used to store user information, transaction details,
    and other relevant data. This project uses the MariaDB version of MySQL,
    which is an open-source fork of MySQL that is fully compatible with it.

    \item \textbf{EJS (Embedded JavaScript):}
    EJS is a simple templating language that lets you generate HTML markup with
    plain JavaScript. It is used to render dynamic content on the server side
    and send it to the client. EJS allows you to include JavaScript code within
    your HTML, making it easy to create dynamic web pages.

    \item \textbf{Git:}
    Git is a distributed version control system that allows you to track
    changes in your codebase. It is widely used for source code management in
    software development. Git allows multiple developers to work on the same
    project simultaneously without conflicts. It provides features like
    branching, merging, and version history, making it easier to collaborate on
    projects.

    \item \textbf{GitHub:}
    GitHub is a web-based platform that uses Git for version control. It
    provides a user-friendly interface for managing Git repositories, making it
    easier to collaborate with other developers. GitHub also offers features
    like issue tracking, pull requests, and project management tools, making it
    a popular choice for open-source projects and team collaboration.

    \item \textbf{Code Editor (Visual Studio Code):}
    Visual Studio Code (VS Code) is a lightweight but powerful source code
    editor that runs on your desktop. It is available for Windows, macOS, and
    Linux. VS Code includes support for debugging, syntax highlighting,
    intelligent code completion, snippets, code refactoring, and embedded
    Git. It also has a rich ecosystem of extensions for additional features
    and functionality.
    \begin{itemize}
        \item Utilize Visual Studio Code as the code editor for project development
        and editing.
    \end{itemize}
\end{enumerate}

\textbf{Installation Steps:}
\begin{enumerate}
    \item \textbf{Install Deno:}
    Follow the official Deno installation guide to set up Deno on your system.
    You can find the installation instructions at
    \url{https://deno.land/#installation}.

    \item \textbf{Install Git:}
    Download and install Git from the official website:
    \url{https://git-scm.com/downloads}. Follow the instructions for your
    operating system.

    \item \textbf{Project Setup:}
    Clone the project repository from GitHub using the following command:
    \begin{lstlisting}[language=bash]
        git clone <repository-url>
    \end{lstlisting}
    Navigate to the project directory:
    \begin{lstlisting}[language=bash]
        cd <project-directory>
    \end{lstlisting}
    Now, simply run the following command to start the server:
    \begin{lstlisting}[language=bash]
        deno run dev
    \end{lstlisting}
    This will install all the required dependencies and start the server.
\end{enumerate}

\chapter{Structure of the Project}

\begin{figure}[h]
    \centering
    % Placeholder for project structure diagram
    \caption{Project Structure}
\end{figure}

\chapter{Interface and Implementation}

\section{app.js}
\begin{lstlisting}[language=JavaScript]
// Code placeholder for app.js
// Contains server setup, database connection, and route handlers
\end{lstlisting}

\section{Account.ejs}
\begin{lstlisting}[language=HTML]
// Code placeholder for Account.ejs
// Contains HTML/EJS template for account creation page
\end{lstlisting}

\section{success.ejs}
\begin{lstlisting}[language=HTML]
// Code placeholder for success.ejs
// Contains HTML/EJS template for success page
\end{lstlisting}

\section{login.ejs}
\begin{lstlisting}[language=HTML]
// Code placeholder for login.ejs
// Contains HTML/EJS template for login page
\end{lstlisting}

\section{login\_success.ejs}
\begin{lstlisting}[language=HTML]
// Code placeholder for login_success.ejs
// Contains HTML/EJS template for login success page
\end{lstlisting}

\section{Transaction.ejs}
\begin{lstlisting}[language=HTML]
// Code placeholder for Transaction.ejs
// Contains HTML/EJS template for transaction page
\end{lstlisting}

\section{Transaction\_details.ejs}
\begin{lstlisting}[language=HTML]
// Code placeholder for Transaction_details.ejs
// Contains HTML/EJS template for transaction details page
\end{lstlisting}

\section{Transaction\_success.ejs}
\begin{lstlisting}[language=HTML]
// Code placeholder for Transaction_success.ejs
// Contains HTML/EJS template for transaction success page
\end{lstlisting}

\section{admin.ejs}
\begin{lstlisting}[language=HTML]
// Code placeholder for admin.ejs
// Contains HTML/EJS template for admin login page
\end{lstlisting}

\section{admin\_success.ejs}
\begin{lstlisting}[language=HTML]
// Code placeholder for admin_success.ejs
// Contains HTML/EJS template for admin success page
\end{lstlisting}

\section{customer\_details.ejs}
\begin{lstlisting}[language=HTML]
// Code placeholder for customer_details.ejs
// Contains HTML/EJS template for customer details page
\end{lstlisting}

\section{transaction\_details.ejs}
\begin{lstlisting}[language=HTML]
// Code placeholder for transaction_details.ejs
// Contains HTML/EJS template for transaction details page
\end{lstlisting}

\section{Header.ejs}
\begin{lstlisting}[language=HTML]
// Code placeholder for Header.ejs
// Contains HTML/EJS template for header component
\end{lstlisting}

\section{Footer.ejs}
\begin{lstlisting}[language=HTML]
// Code placeholder for Footer.ejs
// Contains HTML/EJS template for footer component
\end{lstlisting}

\chapter{Outputs}

\section{Home Page}
\begin{figure}[h]
    \centering
    % Placeholder for home page screenshot
    \caption{Home Page Interface}
\end{figure}

\section{User Sign up}
\begin{figure}[h]
    \centering
    % Placeholder for user signup screenshot
    \caption{User Sign Up Interface}
\end{figure}

\section{User Account Login}
\begin{figure}[h]
    \centering
    % Placeholder for user account login screenshot
    \caption{User Account Login Interface}
\end{figure}

\section{Transaction}
\begin{figure}[h]
    \centering
    % Placeholder for transaction screenshot
    \caption{Transaction Interface}
\end{figure}

\section{Transaction with insufficient money}
\begin{figure}[h]
    \centering
    % Placeholder for transaction with insufficient money screenshot
    \caption{Transaction with Insufficient Money Interface}
\end{figure}

\section{Admin login}
\begin{figure}[h]
    \centering
    % Placeholder for admin login screenshot
    \caption{Admin Login Interface}
\end{figure}

\section{Admin Monitoring Customer Details}
\begin{figure}[h]
    \centering
    % Placeholder for admin monitoring customer details screenshot
    \caption{Admin Monitoring Customer Details Interface}
\end{figure}

\section{Admin Monitoring Customer Transactions}
\begin{figure}[h]
    \centering
    % Placeholder for admin monitoring customer transactions screenshot
    \caption{Admin Monitoring Customer Transactions Interface}
\end{figure}

\chapter{References}

\begin{enumerate}
    \item \raggedright \url{https://stackoverflow.com/questions/5132923/how-to-add-a-primary-key-to-a-mysql-table}
    \item \raggedright \url{https://stackoverflow.com/questions/9070764/insert-auto-increment-primary-key-to-existing-table}
    \item \raggedright \url{https://stackoverflow.com/questions/5813771/in-ejs-template-engine-how-do-i-include-a-footer}
    \item \raggedright \url{https://stackoverflow.com/questions/22212440/using-math-class-to-convert-cosine-into-degrees}
    \item \raggedright \url{https://www.geeksforgeeks.org/how-to-connect-node-js-application-to-mysql/}
    \item \raggedright \url{https://stackoverflow.com/questions/13396119/how-to-add-a-value-to-existing-value-using-sql-query}
\end{enumerate}

\end{document}
\begin{thebibliography}{9}
    \bibitem{nodejs}
    Node.js. \textit{https://nodejs.org/en/}

    \bibitem{express}
    Express.js. \textit{https://expressjs.com/}

    \bibitem{ejs}
    EJS. \textit{https://ejs.co/}

    \bibitem{mysql}
    MySQL. \textit{https://www.mysql.com/}
    \bibitem{visualstudio}
    Visual Studio Code. \textit{https://code.visualstudio.com/}
\end{thebibliography}

\end{document}
